%% feather_iotj.tex
%% FEATHer manuscript — IEEE Internet of Things Journal (IoT-J) submission.
%%
%% Build:
%%   pdflatex feather_iotj
%%   bibtex   feather_iotj
%%   pdflatex feather_iotj
%%   pdflatex feather_iotj
%%
%% Companion files in this directory:
%%   IEEEtran.cls       — journal class (CTAN, v1.8b)
%%   IEEEtran.bst       — bibliography style (CTAN)
%%   feather.bib        — bibliography database (to be populated)
%%   media/             — figures extracted from the original .docx
%%
%% The original TPAMI submission was authored in Word. Body text + math
%% were imported via pandoc from `IEEE_TPAMI_FEATHer_제출본.docx` into
%% `feather_raw.tex` (sibling file) and then transferred section by
%% section into this IEEEtran journal-mode skeleton.

\documentclass[journal]{IEEEtran}

% ---------------- packages ----------------
\usepackage{amsmath,amssymb,amsfonts}
\usepackage{amsthm}
\usepackage{algorithm}
\usepackage{algpseudocode}
\usepackage{booktabs}
\usepackage{multirow}
\usepackage{graphicx}
\graphicspath{{../figures/}{media/media/}}  % renamed figures first, raw fallback
\usepackage{xcolor}
\usepackage{url}
\usepackage{cite}
\usepackage{hyperref}
\hypersetup{
    colorlinks=true,
    linkcolor=blue,
    citecolor=blue,
    urlcolor=blue,
}

% ---------------- theorem environments ----------------
% R8 #4 audit: Section 4 defines Theorems 1-5 here as formal statements;
% Section 6 references them via \ref/\autoref so numbering can never
% drift again.
\newtheorem{theorem}{Theorem}
\newtheorem{lemma}[theorem]{Lemma}
\newtheorem{proposition}{Proposition} % independent counter -> Prop. 1, Thm. 1, 2
\theoremstyle{remark}
\newtheorem{remark}{Remark}

% ---------------- TODO macro ----------------
% Visual placeholder for sentences that depend on Phase 4 / 4b sweep
% data or Cortex-M3 QEMU validation. These show up in the PDF as
% bold red boxes so they cannot be missed during a final read-through.
\usepackage[dvipsnames]{xcolor}
\newcommand{\TODO}[1]{\textcolor{red}{\textbf{[TODO: #1]}}}

% ---------------- title block ----------------
% R2 #1 — old TPAMI title "FEATHer: Fourier-Efficient Adaptive Temporal
% Hierarchy Forecaster for Time-Series Forecasting" repeats "Forecaster
% for Forecasting" and over-promises "Hierarchy". IoT-J reframing
% emphasizes deployment for industrial IoT instead of the method label.
\title{FEATHer: A Deployable Sub-1K Parameter Forecaster for\\
Industrial IoT Edge Devices}

\author{
    Jaehoon Lee,~Seungwoo Lee,~Younghwi Kim,~%
    Dohee Kim,~\IEEEmembership{Member,~IEEE,}~and~%
    Sunghyun Sim,~\IEEEmembership{Member,~IEEE}%
    \thanks{This work was supported by the National Research Foundation
    of Korea (NRF) grant funded by the Korean government (MSIT)
    (No.~RS-2023-00218913) and the Basic Science Research Program
    through the NRF funded by the Ministry of Education
    (No.~RS-2025-25396743).}%
    \thanks{J. Lee, S. Lee, and Y. Kim are with the Graduate School of
    AI Convergence Engineering, Changwon National University, Changwon
    51140, South Korea.}%
    \thanks{D. Kim and S. Sim (corresponding authors) are with the
    Department of AI Convergence, Changwon National University,
    Changwon 51140, South Korea
    (e-mail: \href{mailto:dynamic97312@naver.com}{dynamic97312@naver.com}).}%
}

\markboth{IEEE Internet of Things Journal,~Vol.~XX, No.~X, MONTH~YEAR}%
{Lee \MakeLowercase{\textit{et al.}}: FEATHer — A Deployable Sub-1K
Parameter Forecaster for Industrial IoT Edge Devices}

% ---------------- document ----------------
\begin{document}

\maketitle

% =====================================================================
\begin{abstract}
% R8 #1: "broadly SOTA" claim removed. We report wins honestly and
%        explicitly flag the datasets where larger baselines win.
% R8 #2: sub-1K parameter claim qualified as
%        "configuration-dependent ultra-light regime".
% IoT-J: edge deployment, 4-axis sensor-fault robustness,
%        Wilcoxon significance testing on five seeds.
Time-series forecasting on industrial edge devices --- programmable
logic controllers, embedded microcontrollers --- must satisfy strict
latency, memory, and energy budgets that constrain deployable models to
at most a few thousand parameters. We present FEATHer, a deployable
forecasting system built around four lightweight components: a
multiscale temporal decomposition that splits the input into point-,
high-, mid-, and low-frequency pathways via depthwise 1D convolutions;
a shared Dense Temporal Kernel for temporal mixing without recurrence
or attention; a frequency-aware branch gating that fuses the pathways
using the spectral signature of the normalized input; and a
Sparse-Period Kernel that reconstructs long-horizon outputs through
period-wise reshaping and a shared linear projection. FEATHer operates
in a \emph{configuration-dependent} ultra-light regime --- 453
parameters on ETTh1 and a few thousand on higher-channel datasets ---
and is deployable on Cortex-M3 hardware under $16$--$64$\,KB RAM in
streaming batch-one inference. Across eight long-term benchmarks and
four prediction horizons with five seeds per cell, FEATHer is
competitive in the ultra-light regime against eleven baselines that
span five orders of magnitude in parameter count, including DiPE-Linear
(DASFAA~2026) and MDMLP-EIA (AAAI~2026); larger-capacity baselines
remain stronger on high-channel datasets such as Solar, Traffic, and
Electricity, which we report transparently rather than averaging over.
Beyond clean-input accuracy, we evaluate robustness against four
classes of sensor faults --- additive Gaussian noise, missing values,
impulse outliers, and quantization distortion --- across $19\,200$
corruption conditions, with pairwise Wilcoxon signed-rank tests over
five seeds to support every reported comparison. The results identify a
practical operating regime for next-generation industrial IoT systems
that require continuous, low-power inference under sensor noise and
missing data.
\end{abstract}

\begin{IEEEkeywords}
Time-series forecasting, edge AI, industrial IoT, ultra-lightweight
models, on-device inference, sensor robustness.
\end{IEEEkeywords}

% =====================================================================
\section{Introduction}
\label{sec:introduction}

Time-series forecasting underpins critical functions in modern
industrial systems: production scheduling, predictive maintenance,
energy balancing, anomaly detection, and safety monitoring across
manufacturing, logistics, and transportation infrastructure. As these
domains shift toward cyber-physical and autonomous operation,
forecasting models are increasingly required to run \emph{directly on
edge platforms} --- programmable logic controllers (PLCs), embedded
microcontrollers, industrial IoT sensor nodes, and low-power gateways
--- under strict constraints on latency, memory, and energy. These
platforms typically offer kilobyte-scale weight budgets, single-digit
megahertz cores, and millisecond-level response deadlines, while
operating under nonstationary conditions that include sensor noise,
intermittent communication losses, and quantization-induced
distortion. Large Transformer-based architectures, deep convolutional
encoders, and even moderately sized state-space models are
impractical in this regime; this motivates forecasting designs that
maintain accuracy over long horizons under \emph{extreme} parameter
budgets and that remain reliable under realistic sensor faults.

% R4 #2/3: outdated baselines — both new families included below.
Recent research has explored a range of lightweight forecasting
architectures, including
DLinear~\cite{Zeng2023DLinear}, TiDE~\cite{Das2024TiDE},
TSMixer~\cite{Chen2023TSMixer}, FITS~\cite{Xu2024FITS},
CycleNet~\cite{Lin2024CycleNet},
SparseTSF~\cite{Lin2024SparseTSF}, and very recently
DiPE-Linear~\cite{Wintertee2026DiPELinear} and
MDMLP-EIA~\cite{Zh2026MDMLPEIA}. These methods demonstrate the value
of linear decomposition, shallow temporal mixing, and period-aware
projections. Three challenges nevertheless remain for industrial
long-horizon forecasting. First, many designs rely on a
\emph{single} temporal scale or a fixed periodic structure, which is
insufficient to represent the hierarchical patterns observed in
industrial signals where rapid fluctuations, medium-range
transitions, and long-term seasonal drifts coexist. Second,
lightweight architectures often lack an explicit mechanism for
structured frequency decomposition, forcing heterogeneous temporal
components into a single representational pathway and inducing
cross-frequency interference. Third, despite being labeled
``lightweight,'' many models still operate in the
$10^{3}$--$10^{5}$ parameter range, which can exceed the budgets
imposed by tightly constrained edge controllers
(see Section~\ref{sec:related} for a detailed comparison).

% R8 #2 tone-down: configuration-dependent ultra-light regime.
% R8 #1 tone-down: "competitive" rather than "broadly SOTA".
To address these gaps, we present \textbf{FEATHer}, a deployable
forecasting system designed for accurate long-horizon forecasting
under severe resource constraints. FEATHer follows a structured
design in which representations are explicitly organized across
temporal scales and adaptively fused based on the spectral
characteristics of the input. The model consists of four lightweight
components (Fig.~\ref{fig:overall}): a multiscale temporal
decomposition that splits the instance-normalized input into point-,
high-, mid-, and low-frequency pathways via depthwise 1D
convolutions; a shared Dense Temporal Kernel (DTK) that mixes time
information within each band using a projection-depthwise-projection
sequence, without recurrence or self-attention; a frequency-aware
branch gating module that fuses pathways by computing softmax
weights from the FFT magnitude of the normalized input; and a
Sparse-Period Kernel (SPK) that reconstructs long-horizon outputs
through period-aligned reshaping and a shared linear projection.

In aggregate, these choices place FEATHer in a
\emph{configuration-dependent} ultra-light regime --- 453 parameters
on ETTh1 and at most a few thousand parameters on higher-channel
datasets --- with deployment validated on Cortex-M3 hardware under
$16$--$64$\,KB RAM budgets in streaming batch-one inference
(Section~\ref{sec:edge}). Beyond clean-input accuracy, we evaluate
robustness against four classes of sensor faults that arise in field
deployment --- additive Gaussian noise, missing values, impulse
outliers, and quantization distortion --- across $19\,200$ corruption
conditions, an evaluation axis previously underrepresented in the
LTSF literature
(Section~\ref{sec:robustness}). The contributions of this work are
as follows:

\begin{itemize}
\item[\textbf{(C1)}] An ultra-lightweight \emph{multiscale
  decomposition} that separates input dynamics into four
  frequency-aligned pathways, enabling band-specialized temporal
  modeling under extreme parameter budgets while suppressing
  cross-frequency interference.

\item[\textbf{(C2)}] A shared \emph{Dense Temporal Kernel} (DTK) that
  captures temporal dependencies without recurrence or self-attention
  and prevents parameter growth in multi-branch designs.

\item[\textbf{(C3)}] A \emph{frequency-aware gating} mechanism that
  adaptively fuses the multiscale pathways based on the FFT magnitude
  of the input, providing structured adaptation to nonstationary
  dynamics with negligible parameter overhead.

\item[\textbf{(C4)}] A \emph{Sparse-Period Kernel} (SPK) for
  phase-aligned long-horizon reconstruction that captures periodic
  and seasonal structure with parameter count exactly
  $nm = (L/P)(H/P)$ (Theorem~\ref{thm:spk-minimal}), without
  increasing model depth.

\item[\textbf{(C5)}] A comprehensive benchmark against eleven recent
  baselines that span five orders of magnitude in parameter count
  (from SparseTSF and DiPE-Linear at $\le\!10^{3}$ parameters to
  MDMLP-EIA at $1.6\!\times\!10^{6}$), with paper-default
  per-(method, dataset) hyperparameters, five random seeds, and
  pairwise Wilcoxon signed-rank significance tests on every
  reported comparison (Section~\ref{sec:results}).

\item[\textbf{(C6)}] An IoT-J-targeted \emph{sensor-fault robustness}
  evaluation: $19\,200$ corruption conditions across four fault axes
  and four representative datasets, with the trained checkpoints
  reused at inference time so the result reflects the deployed model
  exactly (Section~\ref{sec:robustness}).

\item[\textbf{(C7)}] A Cortex-M3 \emph{on-device measurement
  protocol} reporting peak RAM, flash usage, activation arena, and
  per-sample latency under $16$--$64$\,KB RAM budgets, making the
  deployability claim verifiable rather than abstract
  (Section~\ref{sec:edge}).
\end{itemize}

The remainder of the paper is organized as follows.
Section~\ref{sec:related} surveys related work on lightweight LTSF
and positions FEATHer with respect to contemporary 2024--2026
baselines. Section~\ref{sec:method} details the four FEATHer
components and the corrected Algorithm~\ref{alg:feather}.
Section~\ref{sec:theory} presents the theoretical guarantees
(DTK Lipschitz stability, SPK parameter minimality, energy-consistent
gating). Section~\ref{sec:setup} states the experimental setup,
Section~\ref{sec:results} the main forecasting results, and
Section~\ref{sec:ablation} the ablation study covering all four
modules across all horizons. Section~\ref{sec:robustness} reports the
four-axis sensor-fault robustness sweep, and
Section~\ref{sec:edge} the on-device measurements.
Section~\ref{sec:conclusion} concludes with limitations and future
directions.

% =====================================================================
% Fig. 1 overall architecture (wide, spans both columns).
% Compiled separately from fig1_overall.tex (TikZ standalone) and
% included here as a PDF inclusion.
\begin{figure*}[t]
    \centering
    \includegraphics[width=0.95\textwidth]{fig1_overall.pdf}
    \caption{Overall architecture of FEATHer. The input window
    $X \in \mathbb{R}^{B \times L \times D}$ passes through (1) the
    multiscale decomposition into point-, high-, mid-, and
    low-frequency pathways; (2) the shared Dense Temporal Kernel (DTK)
    that performs depthwise temporal mixing within each band;
    (3) a frequency-aware branch gating module that derives softmax
    weights from the FFT magnitude of the normalized input; and
    (4) the Sparse-Period Kernel (SPK) that reconstructs the
    long-horizon forecast through period-aligned reshaping with a
    shared linear projection. The four module-detail figures
    (Fig.~\ref{fig:multiscale}--\ref{fig:spk}) expand each component
    in turn.}
    \label{fig:overall}
\end{figure*}

% =====================================================================
\section{Related Work}
\label{sec:related}

Long-term time-series forecasting (LTSF) has been pursued from three
broad angles: (i)~architectures that capture long-range dependencies,
(ii)~efficiency-oriented models that reduce parameter and compute
budgets, and (iii)~methods that introduce explicit temporal structure
through multiscale decomposition, spectral modeling, or periodic
reconstruction. This section reviews each direction, contrasts the
gaps that remain in the ultra-light regime required by industrial
edge deployment, and states which contemporary architectures we
deliberately leave outside the comparison scope.

\subsection{Modeling Long-Horizon Dependencies}
A substantial body of LTSF research scales the Transformer to long
sequences. Autoformer~\cite{Wu2021Autoformer} introduces
decomposition-driven autocorrelation, FEDformer~\cite{Zhou2022FEDformer}
performs attention in the frequency domain,
PatchTST~\cite{Nie2023PatchTST} reformulates the input as channel-
independent patches, and iTransformer~\cite{Liu2024iTransformer}
inverts the attention to operate over variables instead of time steps.
These models achieve strong benchmark accuracy, but their attention
operations and deep stacks typically require $10^{5}$--$10^{6}$
parameters, which is incompatible with the kilobyte-scale weight
budgets of programmable logic controllers and Cortex-class
microcontrollers.

\subsection{Efficiency-Oriented Forecasting}
A parallel line of work pursues competitive accuracy with reduced
budgets. Linear and decomposition-based heads
(DLinear, NLinear~\cite{Zeng2023DLinear}) show that well-structured
linear mappings can match attention-based decoders on standard
benchmarks. Compact MLP-mixing models such as TiDE~\cite{Das2024TiDE}
and TSMixer~\cite{Chen2023TSMixer} push this direction further while
adding shallow nonlinear blocks. Spectral and cycle-aware designs ---
notably FITS~\cite{Xu2024FITS}, which projects the input to a low-rank
frequency representation, and CycleNet~\cite{Lin2024CycleNet}, which
exploits repeating temporal structure --- reduce the effective
parameter count by exploiting compressible temporal regularities.

Two very recent additions extend this lineage and we include both as
contemporary baselines (Section~\ref{sec:results}):
DiPE-Linear~\cite{Wintertee2026DiPELinear} (DASFAA 2026) achieves
strong accuracy with $\sim$$10^{2}$ parameters by combining
disentangled period embeddings with shallow linear heads, and
MDMLP-EIA~\cite{Zh2026MDMLPEIA} (AAAI 2026) reaches sub-2M parameters
through a multi-decomposition MLP with explicit instance adaptation.
Despite this progress, models marketed as ``lightweight'' still
routinely operate in the $10^{3}$--$10^{5}$ parameter range, which is
two to four orders of magnitude above the budgets feasible on
flash-constrained MCUs.

% R1 #6: cite recent hybrid / broad-learning forecasters as
% complementary directions outside the edge-deployment scope.
A complementary line of recent work improves predictive accuracy by
hybridizing broad learning systems or quantum-enhanced recurrent units
with temporal convolution, reporting strong results on chaotic and
financial series~\cite{Su2026BLSAttnTCN, Su2025BLSQLSTM}. These designs
optimize accuracy on server-class hardware rather than the
kilobyte-scale weight budgets that motivate FEATHer, and we note them
as orthogonal advances outside our edge-deployment scope.

\subsection{Structured Multiscale and Spectral Modeling}
A complementary direction adds explicit temporal structure to improve
long-horizon stability. Multi-resolution convolutional designs and
hierarchical decompositions seek to disentangle temporal patterns
through structured transformations. TimesNet~\cite{Wu2023TimesNet}
folds the 1D signal into 2D representations along multiple periods and
applies inception-style convolutions across them.
TimeMixer~\cite{Wang2024TimeMixer} performs multi-scale MLP mixing
along both time and channel axes.
SparseTSF~\cite{Lin2024SparseTSF} introduces period-aware sparse
projections that share linear weights across phases, achieving
remarkable parameter efficiency under explicit periodicity.

FEATHer overlaps with these methods in spirit but differs along three
axes that matter under sub-1K budgets. (i)~Unlike
TimesNet/TimeMixer, the FEATHer multiscale decomposition produces
\emph{exactly four} band-specialized pathways (point, high, mid, low)
with kernel sizes pinned to a fixed receptive-field schedule, removing
the search over multiple periods or hierarchical levels.
(ii)~Frequency-aware gating in FEATHer derives band weights from the
FFT magnitude of the instance-normalized input rather than from
attention or pooled statistics, which keeps the gating module to a
single Conv1D and a softmax. (iii)~The SPK shares the
phase-aligned-projection idea with SparseTSF but applies it
\emph{after} a learnable temporal mixer rather than directly to the
raw signal, which retains parameter minimality
(Theorem~\ref{thm:spk-minimal}) while accommodating non-periodic
fluctuations through the upstream DTK.

\subsection{Architectures Outside the Edge-Deployment Scope}
\label{subsec:scope-exclusion}
Our benchmark deliberately excludes two contemporary families of
LTSF architectures because their deployment profile is incompatible
with our sub-1K, MCU-targeted setting.

\emph{Selective-scan state-space models.}
Mamba-family forecasters such as
S-Mamba~\cite{Wang2024SMamba} achieve linear-time long-range modeling
through the \texttt{selective\_scan\_cuda} and
\texttt{causal\_conv1d\_cuda} kernels. These kernels are CUDA-specific
and have, to our knowledge, no port to ARM Cortex-M or comparable
embedded runtimes; the same model also operates two to three orders of
magnitude above our parameter budget. We therefore cite this family
here but do not include it in the on-device comparison.

\emph{LLM-augmented forecasters.}
TimeCMA~\cite{Liu2025TimeCMA} (AAAI 2025) and related approaches
condition forecasting heads on features extracted from frozen large
language models (e.g., GPT-2, $\sim$$1.2\!\times\!10^{8}$ parameters).
The frozen backbone alone exceeds our parameter budget by five orders
of magnitude and cannot be quantized or pruned away without altering
the method. These approaches lie outside the deployment regime our
work targets and are likewise cited but not benchmarked.

\emph{Foundation forecasting models.}
MOMENT~\cite{Goswami2024Moment} and Moirai~\cite{Woo2024Moirai}
target the universal pre-train-and-transfer regime; both rely on
hundreds-of-millions-parameter Transformer backbones and we exclude
them for the same parameter-budget reason as the LLM-augmented
family above.

\emph{Other concurrent designs we do not benchmark.}
TimeMixer++~\cite{Wang2024TimeMixerPP} extends TimeMixer with a
``pattern machine'' that targets universal predictive analysis; its
default capacity is similar to TimeMixer and inherits the same
$>$$10^{4}$ parameter floor. EDformer~\cite{Chen2024EDformer}
augments the Transformer block with decomposition-aware embeddings
and falls in the same heavy-attention regime as iTransformer and
PatchTST. Both methods are cited here for completeness and target
points in the design space not contested by our sub-1K budget; we
therefore retain the benchmarked subset
(Section~\ref{subsec:baselines}) as the comparison scope.

\subsection{Positioning of FEATHer}
The gap addressed by FEATHer is structural rather than incremental.
High-capacity LTSF models capture long-range dependencies at the cost
of deployability, while existing lightweight models reach a few
thousand parameters but rarely integrate (i)~explicit band-wise
separation, (ii)~spectrum-based adaptive scale selection, and
(iii)~horizon-agnostic period-aligned reconstruction within a single
ultra-compact design. FEATHer combines these three properties in a
$\sim$$10^{2}$--$10^{3}$ parameter envelope deployable on Cortex-M3
hardware under $16$--$64$\,KB RAM budgets
(Section~\ref{sec:edge}), and we evaluate it not only on clean inputs
but also under the four sensor-fault axes documented in
Section~\ref{sec:robustness}.

% =====================================================================
\section{Methodology}
\label{sec:method}

FEATHer is designed to support stable long-horizon forecasting under
the strict latency, memory, and energy constraints typical of
industrial edge devices. The model integrates four lightweight
components illustrated in Fig.~\ref{fig:overall}: a structured
multiscale temporal decomposition, a shared Dense Temporal Kernel
(DTK), a frequency-aware branch gating module, and a Sparse-Period
Kernel (SPK). FEATHer avoids recurrence and self-attention, relying
instead on linear projections, depthwise temporal filtering, and
period-aligned mappings. This section formalizes each component;
the full forward pass is summarized in
Algorithm~\ref{alg:feather}.

\subsection{Problem Setup and Notation}
We consider multivariate forecasting with an input window of length
$L$, a forecasting horizon $H$, and $D$ variables. The input sequence
is $X \in \mathbb{R}^{L \times D}$ and the target horizon is
$Y \in \mathbb{R}^{H \times D}$; FEATHer produces
$\widehat{Y} \in \mathbb{R}^{H \times D}$ from $X$. The model
constructs $B \in \{2,3,4\}$ scale pathways and the active branch
set is $\mathcal{B} \subseteq \{p,h,m,l\}$, corresponding to point,
high, mid, and low frequency pathways with $|\mathcal{B}| = B$. The
DTK uses a latent width $S$, and the SPK uses a period $P$. All
branch representations are aligned to the same temporal length $L$
so that they can be fused without cross-scale alignment overhead.
In this work, ``scale'' refers to frequency-aligned temporal
representations obtained via lightweight temporal filtering.
% R6 #1a: rationale for B in {2,3,4}.
The branch count $B$ is restricted to $\{2,3,4\}$ for two
practical reasons: (i)~the smallest configuration ($B{=}2$, e.g.,
point + low) is sufficient on strongly stationary signals and keeps
the parameter budget minimal; (ii)~$B{=}4$ activates all four
frequency-band specialists and is the default for industrial signals
with mixed-frequency content. We did not observe further gains from
$B \ge 5$ in pilot experiments, and additional branches violate
the sub-1K target on higher-channel datasets.

% R5 #6 + R1 #2: how is the SPK period P chosen?
The SPK period $P$ is treated as an \emph{architectural} hyperparameter
rather than a data-tuned one. In FEATHer we fix $P{=}12$ across all
eight datasets, which satisfies the divisibility constraint
$P\,|\,L$ and $P\,|\,H$ for every horizon $H \in \{96, 192, 336, 720\}$
in the benchmark, and corresponds to a coarse seasonal scale that is
present in most industrial signals (e.g., the 12-hour cycle in
electricity transformer temperature, the 12-step daily cycle at the
1-hour sampling rate, the 12-step diurnal cycle in Weather at the
2-hour resolution). When a deployment scenario has a known dominant
period, $P$ can be set to that value at deployment time without
re-training the rest of the network, since $W \in \mathbb{R}^{n \times m}$
is the only module whose shape depends on $P$.

% R6 #5: notation table promised by Algorithm 1's R8 #5 comment.
Table~\ref{tab:notation} consolidates the symbols used in this
section, in Section~\ref{sec:theory}, and in
Algorithm~\ref{alg:feather}. The two notational conventions
(\,$X^{(b)}, H^{(b)}$ in the body vs.~$h^{(s)}, \widetilde{h}^{(s)}$
in the algorithm) refer to the same tensors at different abstraction
levels; the body uses the branch index $b$ to emphasize parallelism
across $\mathcal{B}$, while the algorithm uses $s$ to make the
sequential pseudocode easier to follow.

\begin{table}[t]
\centering
\caption{FEATHer notation summary. Body and algorithm symbols map
one-to-one through this table (R6 \#5, R8 \#5).}
\label{tab:notation}
\small
\begin{tabular}{ll}
\toprule
Symbol & Meaning \\
\midrule
$X \in \mathbb{R}^{L \times D}$ & input window (length $L$, $D$ vars) \\
$\widetilde X$                  & instance-normalized input \\
$X^{(b)}, H^{(b)}$              & branch input / DTK output (body) \\
$h^{(s)}, \widetilde h^{(s)}$   & same as above (algorithm view) \\
$\mathcal{B} \subseteq \{p,h,m,l\}$ & active branch set \\
$B = |\mathcal{B}| \in \{2,3,4\}$ & number of active branches \\
$S$                              & DTK latent width \\
$P$                              & SPK period \\
$n = L/P,\; m = H/P$             & cycles in input / horizon \\
$W_{\mathrm{in}}, W_{\mathrm{out}}$ & DTK projections ($D{\to}S$, $S{\to}D$) \\
$k_{\mathrm{temp}}, k_{\mathrm{slide}}$ & DTK / SPK depthwise kernels \\
$g \in \Delta^{B}$               & gate weights (softmax over branches) \\
$H_g$ (body) /\;$h$ (alg.)       & gate-fused representation ($L{\times}D$) \\
$H_{\mathrm{agg}}$ /\;$h_{\mathrm{agg}}$ & residual-smoothed fusion \\
$W \in \mathbb{R}^{n \times m}$  & SPK shared linear map \\
$\lambda_{\mathrm{spec}}$        & spectral-separation loss weight \\
\bottomrule
\end{tabular}
\end{table}

\subsection{Structured Multiscale Temporal Decomposition}
\label{subsec:multiscale}
FEATHer first transforms the instance-normalized input $\widetilde X$
into multiple scale representations that emphasize different temporal
behaviors while remaining computationally inexpensive
(Fig.~\ref{fig:multiscale}). For each active branch
$b \in \mathcal{B}$, we generate
$X^{(b)} = \phi^{(b)}(\widetilde X) \in \mathbb{R}^{L \times D}$,
where $\phi^{(b)}(\cdot)$ is a \emph{lightweight depthwise temporal
operator}: a 1D convolution applied independently to each input
variable along time, with the channel dimension preserved. This
construction follows the depthwise convolution pattern popularized by
MobileNet~\cite{Howard2017MobileNets}, restricted here to one
temporal dimension so that the per-branch parameter cost is
$\mathcal{O}(k)$ rather than $\mathcal{O}(D^{2}k)$. The point branch preserves instantaneous
information through a kernel-size-1 depthwise operator, while the
high and mid branches use short-support depthwise convolutions
(kernel sizes $3$ and $5$) to emphasize local and medium-range
variations, respectively. The low branch isolates slow components by
temporally downsampling the input via average pooling with stride
$r$ and then interpolating back to length $L$, yielding an explicit
low-pass bias without introducing heavy parameters.

% R5 #5 defense: asymmetric design rationale.
The asymmetry between the convolution-based and the
pool-then-upsample design choices is deliberate. A learned $k{=}5$
depthwise convolution can model local variations within five
consecutive samples, but its effective receptive field grows only
linearly with kernel size; reaching a low-frequency receptive field
through a single convolution would require kernels of size
$\Theta(L/r)$, which dominates the parameter budget on
high-channel datasets. The pool-then-upsample chain in the Low
branch instead obtains an effective receptive field of $r$ samples
\emph{without} learnable parameters in the pool stage, which leaves
budget for the shared linear projection that follows. This
parameter-receptive-field trade-off is the structural reason the
four branches use the operators they do.

\subsection{Dense Temporal Kernel (DTK)}
Each decomposed pathway is processed by a shared DTK that performs
efficient temporal mixing without recurrence or self-attention
(Fig.~\ref{fig:dtk}). DTK adopts a
projection--depthwise-filtering--inverse-projection structure. For a
branch input $X^{(b)} \in \mathbb{R}^{L \times D}$, DTK first
projects the input to a compact latent width $S$:
\begin{equation}
\label{eq:dtk-zin}
Z^{(b)} = X^{(b)}\, W_{\mathrm{in}} \in \mathbb{R}^{L \times S},
\end{equation}
with $W_{\mathrm{in}} \in \mathbb{R}^{D \times S}$. A depthwise
temporal convolution then mixes information along time
\emph{independently per latent channel} using the shared kernel
$k_{\mathrm{temp}}$:
\begin{equation}
\label{eq:dtk-dwconv}
U^{(b)} = \mathrm{DWConv}_t(Z^{(b)};\, k_{\mathrm{temp}})
       \in \mathbb{R}^{L \times S}.
\end{equation}
The latent representation is then projected back to the model
dimension:
\begin{equation}
\label{eq:dtk-hout}
H^{(b)} = U^{(b)}\, W_{\mathrm{out}} \in \mathbb{R}^{L \times D},
\end{equation}
with $W_{\mathrm{out}} \in \mathbb{R}^{S \times D}$. The DTK
parameters $\{W_{\mathrm{in}}, W_{\mathrm{out}}, k_{\mathrm{temp}}\}$
are shared across all branches $b \in \mathcal{B}$. This sharing
prevents parameter growth when the branch count grows from $B{=}2$
to $B{=}4$ while still allowing each branch to express
scale-specific dynamics through its distinct input $X^{(b)}$. The
Lipschitz stability of this construction is established in
Theorem~\ref{thm:dtk-lipschitz}.

\subsection{Frequency-aware Branch Gating}
\label{subsec:gate}
The importance of each temporal scale varies across instances and
operating regimes. FEATHer therefore employs a frequency-aware
gating mechanism that assigns branch weights from the input spectral
profile (Fig.~\ref{fig:gate}). Starting from the instance-normalized
input $\widetilde X$, we compute the real FFT along the temporal
axis,
\begin{equation}
\label{eq:gate-fft}
F = \mathcal{F}(\widetilde X), \qquad A = |F|,
\end{equation}
and obtain a compact spectral descriptor by averaging the magnitude
across variables,
\begin{equation}
\label{eq:gate-meanD}
a = \frac{1}{D} \sum_{d=1}^{D} A_{:,d} \in \mathbb{R}^{L_f},
\quad L_f = \lfloor L/2 \rfloor + 1.
\end{equation}
A lightweight gating network $\psi(\cdot)$ maps $a$ to branch
logits,
\begin{equation}
\label{eq:gate-psi}
u = \psi(a) \in \mathbb{R}^{B},
\end{equation}
and a softmax produces non-negative weights,
\begin{equation}
\label{eq:gate-softmax}
g = \mathrm{softmax}(u) \in \Delta^{B}.
\end{equation}
The fused representation is then
\begin{equation}
\label{eq:gate-fuse}
H_g = \sum_{b \in \mathcal{B}} g_b\, H^{(b)} \in \mathbb{R}^{L \times D}.
\end{equation}

% R6 #1d: dimension transition L_f -> B inside psi(.).
The transition from
$a \in \mathbb{R}^{L_f}$ in~\eqref{eq:gate-meanD} to
$u \in \mathbb{R}^{B}$ in~\eqref{eq:gate-psi} occurs inside
$\psi(\cdot)$, which we implement as a small one-dimensional
convolution stack followed by global average pooling and a final
linear projection to $B$: explicitly, $\psi(a) = W_g\,
\mathrm{Pool}\!\bigl(\mathrm{Conv1D}(a)\bigr)$ with
$W_g \in \mathbb{R}^{C_g \times B}$ for a small intermediate channel
width $C_g$. This composition reduces a length-$L_f$ descriptor to
$B$ logits while keeping the gate module under a hundred parameters.
The energy-consistent interpretation of $g$ is stated formally in
Proposition~\ref{prop:energy-gate}.

\subsection{Sparse-Period Kernel (SPK) for Long-Horizon Reconstruction}
\label{subsec:spk}
To produce $\widehat Y$ efficiently, FEATHer employs a Sparse-Period
Kernel that reconstructs periodic and seasonal structure through
compact period-aligned transformations (Fig.~\ref{fig:spk}). Starting
from $H_g \in \mathbb{R}^{L \times D}$, SPK applies a depthwise
residual aggregation that smooths short-term fluctuations while
preserving the input length $L$:
\begin{equation}
\label{eq:spk-residual}
H_{\mathrm{agg}} = H_g + \mathrm{DWConv}_t(H_g;\, k_{\mathrm{slide}}).
\end{equation}
SPK then reorganizes $H_{\mathrm{agg}}$ into phase-aligned groups using
a period $P$. When $L$ and $H$ are divisible by $P$, we define
$n = L/P$ and $m = H/P$. For each variable $d$, the channel signal
$(H_{\mathrm{agg}})_{:,d}$ is
reshaped into a matrix $\widetilde H_d \in \mathbb{R}^{P \times n}$
whose rows correspond to fixed phases within the period. A
\emph{single} shared linear map $W \in \mathbb{R}^{n \times m}$ is
applied along the period axis:
\begin{equation}
\label{eq:spk-linear}
\widehat Y_d(p, :) = \widetilde H_d(p, :)\, W,
\qquad p = 1, \ldots, P.
\end{equation}
The phase-wise predictions are reassembled by interleaving phases
to form $\widehat Y \in \mathbb{R}^{H \times D}$. When $L$ or $H$ is
not divisible by $P$, we pad to the next multiple of $P$ and crop
the final output back to length $H$. Sharing $W$ across phases
\emph{and} variables yields a compact horizon mapping with exactly
$nm = (L/P)(H/P)$ scalar parameters, which matches the parameter
lower bound established in Theorem~\ref{thm:spk-minimal}.

\subsection{End-to-End FEATHer Pipeline}
Given $X$, FEATHer first generates the multiscale signals
$\{X^{(b)}\}_{b \in \mathcal{B}}$, processes each pathway through the
shared DTK to obtain $\{H^{(b)}\}_{b \in \mathcal{B}}$, computes
spectrum-conditioned weights $g$, and fuses the branch
representations into $H_g$. The SPK then produces the final forecast
$\widehat Y$. The model is trained end-to-end by minimizing
$\mathcal{L}(\widehat Y, Y) = \mathrm{MSE}(\widehat Y, Y) +
\lambda_{\mathrm{spec}} \mathcal{L}_{\mathrm{spec}}$, where
$\mathcal{L}_{\mathrm{spec}}$ is the spectral-separation auxiliary
loss specific to FEATHer
% R8 #5 last bullet: training-objective ambiguity removed.
(the precise form of $\mathcal{L}_{\mathrm{spec}}$ is given in the
appendix of the released code). The forward computation is identical
during training and inference; training adds the gradient update
shown in line~\ref{alg:feather}-29 of Algorithm~\ref{alg:feather}.

% R1 #3: bridge from Sec III modules to Sec IV guarantees.
Three of the four components carry verifiable guarantees that we
return to in Section~\ref{sec:theory}. The DTK is the only block
that introduces a learnable temporal mixer; we therefore bound its
input-to-output Lipschitz constant in
Theorem~\ref{thm:dtk-lipschitz} so that sensor-noise robustness
(Section~\ref{sec:robustness}) is governed by an explicit, small set
of weights. The SPK head implements a phase-shared linear map and we
show in Theorem~\ref{thm:spk-minimal} that its $nm$ parameter
count is the lower bound for that function class --- the SPK is not
under-parameterized relative to what it can represent. The
frequency-aware gate is justified through
Proposition~\ref{prop:energy-gate} as an entropy-regularized energy
selector; the proposition does not claim optimality and we make
that scope explicit in
Remark~\ref{rem:gate-justify}. The multiscale decomposition is
parameter-light by construction and does not require a separate
guarantee.

% --- Module-detail figures (Fig 2-5) included from standalone TikZ files ---
\begin{figure*}[t]
    \centering
    \includegraphics[width=0.85\textwidth]{fig2_multiscale.pdf}
    \caption{Multiscale Decomposition (FEATHer module \textbf{[1]},
    Fig.~\ref{fig:overall}). Four parallel branches separate the
    instance-normalized input $\widetilde{X}$ into point ($k{=}1$),
    high ($k{=}3$), mid ($k{=}5$), and low (AvgPool $\to$ Linear
    $\to$ Interp) frequency bands. The asymmetric Low branch obtains
    a much coarser receptive field than a single $k{=}5$ convolution
    could provide.}
    \label{fig:multiscale}
\end{figure*}

\begin{figure*}[t]
    \centering
    \includegraphics[width=0.85\textwidth]{fig3_dtk.pdf}
    \caption{Shared Dense Temporal Kernel (DTK, FEATHer module
    \textbf{[2]}). The same DTK weights are reused across the four
    frequency branches. Inside the DWConv block, $S$ parallel lanes
    share the kernel $k_{\text{temp}}$ but are processed
    \emph{independently per latent channel}: channel $s$'s output
    uses only channel $s$'s input, with no cross-channel mixing.}
    \label{fig:dtk}
\end{figure*}

\begin{figure*}[t]
    \centering
    \includegraphics[width=0.95\textwidth]{fig4_gate.pdf}
    \caption{Frequency-Aware Branch Gating (FEATHer module
    \textbf{[3]}). The gate weights are computed in seven steps:
    FFT magnitude, mean over channels $D$, Conv1D into 4 bands, mean
    over frequency bins $L_f$, and softmax, yielding a probability
    vector $g \in \mathbb{R}^{4}$ that fuses the four DTK outputs.}
    \label{fig:gate}
\end{figure*}

\begin{figure*}[t]
    \centering
    \includegraphics[width=0.95\textwidth]{fig5_spk.pdf}
    \caption{Sparse-Period Kernel (SPK, FEATHer module \textbf{[4]}).
    After a residual DWConv, the signal is reshaped into a
    $P{\times}n$ grid of $P$ phases and $n{=}L/P$ periods. A
    \emph{single} shared matrix $W \in \mathbb{R}^{n \times m}$ maps
    each phase row to the corresponding $m{=}H/P$ output positions
    via Eq.~\eqref{eq:spk-linear}. The result is interleaved back
    into a one-dimensional forecast $\widehat{Y}$.}
    \label{fig:spk}
\end{figure*}

% --- Algorithm 1 — corrected per R8 #5 audit (drafts/05_algorithms_word.md) ---
% R8 #5 fixes applied inline (flagged in comments):
%   (a) line for the high branch: h^(p) was used in the TPAMI version,
%       corrected here to h^(h).
%   (b) "MeanChannel(M))" had a stray ')'; corrected to MeanChannel(M).
%   (c) Body uses X^(b)/H^(b), algorithm uses h^(s)/h-tilde^(s); the
%       notation table in the body Sec 3 maps the two.
\begin{algorithm}[t]
\caption{FEATHer training step.}
\label{alg:feather}
\begin{algorithmic}[1]
\Require Input window $X \in \mathbb{R}^{L \times D}$;
         branch count $B \in \{2,3,4\}$, active set
         $\mathcal{B} \subseteq \{p,h,m,l\}$;
         SPK period $P$; horizon $H$; parameters $\Theta$.
\Ensure  Forecast $\widehat{Y} \in \mathbb{R}^{H \times D}$.

\Statex \emph{Step 1 -- Structured multiscale temporal decomposition.}
\State $\widetilde{X} \leftarrow \mathrm{InstanceNorm}(X)$
\State $h^{(p)} \leftarrow \mathrm{DWConv}_{1\mathrm{D}}(\widetilde{X};\, k{=}1)$
        \Comment{point branch}
\If{$B = 4$}
    \State $h^{(h)} \leftarrow \mathrm{DWConv}_{1\mathrm{D}}(\widetilde{X};\, k{=}3)$
        \Comment{high branch \textbf{(R8 \#5: was $h^{(p)}$ in TPAMI)}}
\EndIf
\If{$B \ge 3$}
    \State $h^{(m)} \leftarrow \mathrm{DWConv}_{1\mathrm{D}}(\widetilde{X};\, k{=}5)$
        \Comment{mid branch}
\EndIf
\State $x_{\downarrow} \leftarrow \mathrm{AvgPool}_{1\mathrm{D}}(\widetilde{X};\, r{=}4)$
\State $h^{(\ell)} \leftarrow \mathrm{Upsample}_{1\mathrm{D}}(x_{\downarrow};\, L)$
        \Comment{low branch}

\Statex \emph{Step 2 -- Shared Dense Temporal Kernel (DTK).}
\For{each active branch $s \in \mathcal{B}$}
    \State $z^{(s)} \leftarrow h^{(s)}\, W_{\mathrm{in}}$
    \State $u^{(s)} \leftarrow \mathrm{DWConv}(z^{(s)};\, k_{\mathrm{temp}})$
    \State $\widetilde{h}^{(s)} \leftarrow u^{(s)}\, W_{\mathrm{out}}$
\EndFor

\Statex \emph{Step 3 -- Frequency-aware branch gating.}
\State $F \leftarrow \mathcal{F}(\widetilde{X})$,\quad $M \leftarrow |F|$
\State $a \leftarrow \mathrm{MeanChannel}(M)$
        \Comment{\textbf{R8 \#5: removed stray ')' from TPAMI ``MeanChannel(M))''}}
\State $z_g \leftarrow \mathrm{Conv1D}(a)$
\State $g \leftarrow \mathrm{softmax}(\mathrm{Pool}(z_g))$
        \Comment{$g \in \mathbb{R}^{B}$, $\sum_s g_s = 1$}
\State $h \leftarrow \sum_{s \in \mathcal{B}} g_s \, \widetilde{h}^{(s)}$

\Statex \emph{Step 4 -- Sparse-Period Kernel (SPK).}
\State $h_{\mathrm{agg}} \leftarrow h + \mathrm{DWConv}(h;\, k_{\mathrm{slide}})$
\State Reshape each channel of $h_{\mathrm{agg}}$ into $\widetilde{H}_d \in \mathbb{R}^{P \times n}$,
       where $n = L/P$.
\For{each channel $d = 1,\ldots,D$ and phase $p = 1,\ldots,P$}
    \State $Y_d(p, :) \leftarrow \widetilde{H}_d(p, :)\, W$
        \Comment{$W \in \mathbb{R}^{n \times m}$, $m = H/P$}
\EndFor
\State Reassemble $\{Y_d\}_d$ into $\widehat{Y} \in \mathbb{R}^{H \times D}$
       by phase interleaving.

\Statex \emph{Step 5 -- Loss and parameter update.}
\State $L \leftarrow \mathrm{MSE}(\widehat{Y}, Y) + \lambda_{\mathrm{spec}} L_{\mathrm{spec}}$
\State $\Theta \leftarrow \Theta - \eta\, \nabla_\Theta L$
\end{algorithmic}
\end{algorithm}

% =====================================================================
\section{Theoretical Analysis}
\label{sec:theory}

% R8 #4 resolution: TPAMI body referenced Theorems 3/4/5 in Sections
% 6.3-6.5 but only Theorems 1 and 2 were ever stated. Rather than
% manufacture three new formal claims (which R2 #6/#7, R4 #4, R7 #2
% already flagged as over-claiming relative to substance), this revision
% retains only the two verifiable theorems plus a single proposition
% covering the gate, and rewrites every Sec 6.x reference to point to
% \ref{thm:dtk-lipschitz}, \ref{thm:spk-minimal}, \ref{prop:energy-gate},
% or the empirical results in Sec 5.

This section provides theoretical support for FEATHer by formalizing
(i)~the stability of the shared DTK, (ii)~the parameter minimality of
the period-aligned reconstruction in the SPK, and (iii)~an
energy-consistent interpretation of the spectral gate. We deliberately
restrict the analysis to guarantees that correspond directly to the
implemented operators and that remain verifiable under sub-1K parameter
budgets; stronger claims (e.g., strict spectral disjointness) would
require assumptions that are not necessary to support FEATHer's
contributions.

\subsection{Preliminaries and Operator Notation}

We view each temporal module in FEATHer as a linear operator acting
along the temporal dimension and applied channel-wise unless stated
otherwise. For $X \in \mathbb{R}^{L \times D}$, we denote the induced
operator norm by $\|\cdot\|_2$. Let $\mathrm{DWConv}_t(\cdot;k)$ be a
depthwise 1D convolution along time. In implementation, the DTK
temporal convolution uses left padding followed by output trimming to
preserve sequence length $L$, yielding a causal-style local mixing
operator. The SPK aggregation uses a residual form
$H + \mathrm{DWConv}_t(H; k_{\mathrm{slide}})$.

\begin{remark}[Implementation faithfulness]
\label{rem:impl-faithful}
All statements below are written for the operators \emph{as
implemented} (left padding + trimming in DTK; residual aggregation in
SPK). This prevents theory-implementation mismatch and makes the
analysis directly reproducible from the released code.
\end{remark}

\subsection{Stability of the DTK}
DTK applies a linear projection, depthwise temporal filtering, and an
inverse projection:
\begin{equation}
\label{eq:dtk}
\mathrm{DTK}(H) = \mathrm{DWConv}_t(H\, W_{\mathrm{in}};\, k_{\mathrm{temp}})\,
                  W_{\mathrm{out}},
\end{equation}
where $W_{\mathrm{in}} \in \mathbb{R}^{D \times S}$ and
$W_{\mathrm{out}} \in \mathbb{R}^{S \times D}$.

\begin{theorem}[Global Lipschitz stability of the shared DTK]
\label{thm:dtk-lipschitz}
Let $\mathcal{K}$ denote the linear operator induced by the implemented
depthwise temporal convolution in DTK (left padding + trimming). Then,
for any $H, H' \in \mathbb{R}^{L \times D}$,
\begin{equation*}
\|\mathrm{DTK}(H) - \mathrm{DTK}(H')\|_2 \;\le\;
  \|W_{\mathrm{in}}\|_2 \, \|\mathcal{K}\|_2 \, \|W_{\mathrm{out}}\|_2 \,
  \|H - H'\|_2.
\end{equation*}
Moreover, if each depthwise kernel satisfies
$\|k_{\mathrm{temp}}\|_1 \le \kappa$, then
$\|\mathcal{K}\|_2 \le \kappa$, so DTK is globally Lipschitz with
constant $\|W_{\mathrm{in}}\|_2\, \kappa\, \|W_{\mathrm{out}}\|_2$.
\end{theorem}

\begin{proof}
DTK is a composition of linear operators, and the induced-norm bound
follows from submultiplicativity. For depthwise convolution, the
operator norm is bounded by the $\ell_1$ norm of the kernel (standard
discrete-time filtering bound). Left padding and trimming preserve
sequence length and do not increase the induced operator norm beyond
that of the corresponding convolution.
\end{proof}

\begin{remark}[Implications for long-horizon forecasting]
\label{rem:dtk-impl}
Theorem~\ref{thm:dtk-lipschitz} guarantees that the shared temporal
mixer cannot arbitrarily amplify perturbations in the branch
representations. This matters for industrial signals where sensor
noise and operating-regime shifts may otherwise destabilize
extrapolation. The bound also clarifies that stability is governed by a
small set of parameters
($W_{\mathrm{in}}, W_{\mathrm{out}}, k_{\mathrm{temp}}$), which is
desirable in ultra-compact regimes.
\end{remark}

\subsection{Energy-consistent Interpretation of the Frequency Gate}
FEATHer computes branch weights from the magnitude spectrum of the
instance-normalized input. Let $\widetilde{X}$ be that input and let
$a$ be a reduced spectral descriptor (e.g., channel-averaged
magnitude). A lightweight gating network produces logits and weights
$g \in \Delta^B$ (the $B$-simplex), and FEATHer fuses the branch
outputs as $H = \sum_{s \in \mathcal{B}} g_s\, H^{(s)}$.

\begin{proposition}[Energy-consistent soft selection]
\label{prop:energy-gate}
Let $E_s \ge 0$ denote a branch relevance score derived from the
spectral energy captured by branch $s$. The entropy-regularized
surrogate
\begin{equation*}
\max_{g \in \Delta^B} \;\sum_{s \in \mathcal{B}} g_s E_s
       \;+\; \tau \sum_{s \in \mathcal{B}} g_s \log g_s
\end{equation*}
is solved by $g_s \propto \exp(E_s / \tau)$, a soft selection rule that
is monotone in $E_s$ and approaches hard selection as $\tau \to 0$.
\end{proposition}

\begin{remark}[Why we use spectrum]
\label{rem:gate-justify}
Proposition~\ref{prop:energy-gate} gives a principled justification for
the design choice \emph{without} overclaiming optimality. Using
spectral features to compute $g$ is consistent with selecting the most
relevant temporal scales for each instance, while the learned gating
network keeps parameter overhead negligible.
\end{remark}

\subsection{Parameter Minimality of the SPK}
SPK first applies a residual aggregation
$H_{\mathrm{agg}} = H + \mathrm{DWConv}_t(H; k_{\mathrm{slide}})$,
reshapes each channel into a phase-aligned representation with period
$P$, and applies a shared linear map $W \in \mathbb{R}^{n \times m}$,
where $n = L/P$ and $m = H/P$.

\begin{theorem}[Parameter minimality of the SPK linear head]
\label{thm:spk-minimal}
Consider the class of functions that map, for each phase, the $n$
observed cycle values to $m$ future cycle values using a
phase-\emph{shared} linear map. Any linear operator realizing this
mapping requires at least $nm$ scalar degrees of freedom. The SPK head
uses exactly $nm$ parameters in $W$ and is therefore
parameter-minimal within this class.
\end{theorem}

\begin{proof}
A linear map from $\mathbb{R}^n$ to $\mathbb{R}^m$ has $nm$ degrees of
freedom. SPK applies such a map in the phase-aligned space using a
shared $W$, exactly matching this lower bound.
\end{proof}

\begin{remark}[Scope of the guarantee]
\label{rem:spk-scope}
Theorem~\ref{thm:spk-minimal} does \emph{not} claim that all time
series are periodic. It states that when long-horizon structure is
well explained by phase-aligned cycle dynamics, SPK allocates
parameters exactly to the intrinsic degrees of freedom of that
structure, thereby avoiding horizon-specific decoders. This inductive
bias is particularly effective under sub-1K parameter budgets.
\end{remark}

\subsection{Complexity and the Sub-1K Regime}
DTK parameters are $2DS + S\, k_{\mathrm{temp}}$; SPK parameters are
$k_{\mathrm{slide}} + nm$. The decomposition filters and the gating
module introduce only minor overhead, and because DTK is shared across
branches, the parameter count does not grow with $|\mathcal{B}|$.
Runtime complexity is linear in $L$ and $D$, up to small constants
determined by $k_{\mathrm{temp}}$, $k_{\mathrm{slide}}$, and the
latent width $S$.

\begin{remark}[Why we avoid stronger theoretical claims]
\label{rem:scope}
Under compact designs, the most meaningful guarantees are those tied
to the implemented operators and to measurable quantities: stability
bounds, parameter minimality, and scaling behavior. Stronger claims
such as strict spectral disjointness or near-orthogonality would
require assumptions difficult to verify in practice and are not
necessary to support FEATHer's core contributions.
\end{remark}

% =====================================================================
\section{Experimental Setup}
\label{sec:setup}

This section states the experimental protocol that governs every
quantitative result reported in Sections~\ref{sec:results}--%
\ref{sec:robustness}. The same protocol is used across all baselines
and FEATHer to keep the comparison fair, and is implemented
end-to-end in the released code so that every cell in the result
tables is reproducible from a single command.

\subsection{Datasets}
\label{subsec:datasets}
We evaluate on eight widely used multivariate long-term forecasting
benchmarks introduced or popularized by the Autoformer
benchmark~\cite{Wu2021Autoformer}:
\textbf{ETTh1, ETTh2, ETTm1, ETTm2} (electricity transformer
temperature, $D{=}7$ channels each), \textbf{Weather}
($D{=}21$ atmospheric channels, 10-minute sampling),
\textbf{Exchange} (daily exchange rates of $D{=}8$ currencies),
\textbf{Electricity} (hourly consumption from $D{=}321$ clients,
UCI repository), and \textbf{Traffic} (hourly occupancy from
$D{=}862$ San Francisco freeway sensors, California PEMS portal).
The datasets cover sampling resolutions from 10 minutes to one day,
dimensionality from $D{=}7$ to $D{=}862$, and durations from
$\sim$$10^{4}$ to $\sim$$10^{5}$ time steps. All datasets are
processed in strict chronological order; the train, validation, and
test splits use the canonical $6{:}2{:}2$ ratio for the ETT family
and $7{:}1{:}2$ for the remaining datasets, matching the conventions
of prior work
(e.g.,~\cite{Wu2021Autoformer,Nie2023PatchTST,Liu2024iTransformer}).
Each input variable is standardized using $z$-score normalization
computed on the training split only; the same normalization
parameters are applied to the validation and test splits.

% R6 #6a: variable selection on Traffic.
On Traffic, we keep \emph{all} $862$ sensor channels rather than
selecting a subset, matching the convention used by
PatchTST~\cite{Nie2023PatchTST} and
iTransformer~\cite{Liu2024iTransformer}; this keeps the comparison
directly aligned with the published numbers and also exercises every
baseline at its highest-dimensionality cell, which surfaces the
peak-RAM scaling differences reported in Section~\ref{sec:edge}.

% R1 #1: data preprocessing impact.
We deliberately do \emph{not} tune preprocessing per dataset; the
same normalization, same split ratios, and same input length
$L{=}96$ are used across all eight datasets and for every baseline.
This isolates architectural differences from preprocessing
artefacts and lets us interpret the reported metrics as purely
model-driven, at the cost of leaving room for preprocessing-aware
gains that we leave to future work.

For each dataset we evaluate four forecasting horizons
$H \in \{96, 192, 336, 720\}$ time steps with a fixed input length
$L = 96$, so that every model produces the entire horizon in a
single forward pass without autoregressive decoding.

\subsection{Baselines}
\label{subsec:baselines}
We compare FEATHer against eleven contemporary baselines that span
five orders of magnitude in parameter count (Table~\ref{tab:params}):
SparseTSF~\cite{Lin2024SparseTSF} (ICML 2024 oral), DiPE-Linear
\cite{Wintertee2026DiPELinear} (DASFAA 2026),
FITS~\cite{Xu2024FITS} (ICLR 2024 spotlight),
DLinear~\cite{Zeng2023DLinear} (AAAI 2023),
PatchTST~\cite{Nie2023PatchTST} (ICLR 2023),
TimeMixer~\cite{Wang2024TimeMixer} (ICLR 2024),
TQNet~\cite{Lin2025TQNet} (ICML 2025),
LMS-AutoTSF~\cite{Ibrahim2024LMSAutoTSF}, TimesNet~\cite{Wu2023TimesNet}
(ICLR 2023), iTransformer~\cite{Liu2024iTransformer} (ICLR 2024), and
MDMLP-EIA~\cite{Zh2026MDMLPEIA} (AAAI 2026). The two 2026 baselines
(DiPE-Linear and MDMLP-EIA) directly address the
``outdated-baselines'' concern raised against the TPAMI submission of
this work. Architectures outside the deployment scope of this paper
(S-Mamba and TimeCMA) are discussed in
Section~\ref{subsec:scope-exclusion} but not benchmarked here.

\begin{table}[t]
\centering
\caption{Parameter counts on ETTh1 ($D{=}7$) with the paper-default
configuration of each method, sorted by size.}
\label{tab:params}
\small
\begin{tabular}{lrr}
\toprule
Method            & Venue         & Params (ETTh1) \\
\midrule
DiPE-Linear       & DASFAA 2026   & 267            \\
SparseTSF         & ICML 2024     & 41             \\
\textbf{FEATHer (ours)} & ---     & \textbf{453}   \\
FITS              & ICLR 2024     & 1{,}200        \\
DLinear           & AAAI 2023     & 18{,}624       \\
PatchTST          & ICLR 2023     & 35{,}171       \\
TimeMixer         & ICLR 2024     & 75{,}497       \\
TQNet             & ICML 2025     & 86{,}192       \\
LMS-AutoTSF       & 2024          & 181{,}318      \\
TimesNet          & ICLR 2023     & 605{,}479      \\
iTransformer      & ICLR 2024     & 841{,}568      \\
MDMLP-EIA         & AAAI 2026     & 1{,}632{,}832  \\
\bottomrule
\end{tabular}
\end{table}

\subsection{Training Protocol}
\label{subsec:protocol}

% R8 #3 defense: explicit per-(method, dataset) HP table + code repo cite.
We follow a TFB-style protocol in which a small set of axes is
unified across all (method, dataset) cells, while
architecture-sensitive hyperparameters are taken from each baseline's
\emph{official per-dataset script} rather than re-tuned on this
benchmark.

\paragraph{Unified axes (every cell).} Input length $L = 96$;
batch size $32$; training epochs $50$; optimizer
\textsc{Adam} with cosine-annealing learning-rate schedule; cuDNN
deterministic mode (\texttt{torch.backends.cudnn.deterministic =
True}); automatic mixed precision (AMP) \emph{disabled} for fairness
and to avoid numerical artefacts in FITS (complex-valued gradients)
and TimesNet (half-precision cuFFT on non-power-of-two sizes); five
seeds $\{2025, 2026, 2027, 2028, 2029\}$ per cell.
% R2 #16: epoch budget + early stopping.
Within the $50$-epoch budget we apply early stopping on the
validation loss with a uniform patience of $10$ epochs. Model
selection follows the protocol used by every baseline's official
implementation: the checkpoint of the best-validation epoch is
retained, and the test split is evaluated \emph{exactly once} on
that selected model; the test set never participates in model
selection, hyperparameter choice, or stopping decisions. The same
selected checkpoint is reused verbatim for the robustness
evaluation of Section~\ref{sec:robustness}.

\paragraph{Per-(method, dataset) axes.}
Learning rate, training loss (MSE for baselines; $L_1 + \lambda
L_{\mathrm{spec}}$ for FEATHer), architecture widths
($d_{\mathrm{model}}, d_{\mathrm{ff}}, e_{\mathrm{layers}},
n_{\mathrm{heads}}$, dropout), and any dataset-frequency-tied
constants (e.g.,~the cycle period of TQNet, the period of SparseTSF)
are read from each baseline repository's \textit{official} script
for the corresponding ETT / Weather / Electricity / Traffic / Exchange
configuration. The complete set of overrides is shipped as a Python
table in the released code (\texttt{baselines/\_\_init\_\_.py},
fields \texttt{\_METHOD\_DEFAULTS} and \texttt{\_DATASET\_OVERRIDES})
and reproduced in the appendix.

\paragraph{FEATHer's protocol --- per-dataset selection.}
Like every baseline, \emph{FEATHer's hyperparameters are selected
per dataset on the validation split}. Starting from a canonical
configuration ($S{=}8$, $k_{\mathrm{temp}}{=}7$, $P{=}12$,
$\lambda_{\mathrm{spec}}{=}0.01$, $\mathrm{lr}{=}10^{-3}$,
MAE\,$+\,\lambda_{\mathrm{spec}}$ spectral-separation loss), a
one-factor-at-a-time search over the validation set fixes the final
per-dataset values, which are shipped in the same
\texttt{\_DATASET\_OVERRIDES} table as the baselines. This ensures
an apples-to-apples comparison: every method, FEATHer included, is
evaluated at its own best validation operating point rather than a
hand-picked configuration.

\subsection{Evaluation Protocol}
\label{subsec:evaluation}
We report mean squared error (MSE) and mean absolute error (MAE) on
the test split for every (model, dataset, horizon) cell, averaged
over the five seeds. Standard deviations are reported in the
appendix tables.
% R6 #6c: motivation for MSE/MAE only (no COR/R^2 in the main table).
We deliberately use only MSE and MAE in the main results table:
both have a clear lower-is-better orientation that avoids the
mixed-direction caption pitfall raised against the TPAMI submission
(R6 \#8), and both directly track the squared- and absolute-loss
objectives optimized during training. Correlation-based metrics
(Pearson $r$ or $R^{2}$) are reported only in the per-seed appendix
because their interpretation differs from MSE/MAE and can mislead at
short horizons when the conditional mean is approximately constant.

To assess statistical significance, we apply the non-parametric
\emph{Wilcoxon signed-rank test} on the pairwise per-seed score
differences between FEATHer and each baseline, following the
convention adopted in the recent TFB and CF-JEPA benchmarks.
A pair is flagged as significant at $p < 0.05$ (two-sided).

\subsection{Hardware}
\label{subsec:hardware}
Training and evaluation are performed on a server with an NVIDIA
GPU and PyTorch~\mbox{2.5}. On-device latency and memory
measurements (Section~\ref{sec:edge}) are obtained on a physical
LM3S6965EVB development board (ARM Cortex-M3, 50\,MHz,
64\,KB on-chip SRAM), compiled with
\texttt{arm-none-eabi-gcc}~12.2.1 using the standard
\texttt{-Os -mthumb -mcpu=cortex-m3} flags. Three RAM budgets
$\{16, 32, 64\}$\,KB are enforced through the linker script so that
the inference runtime cannot exceed the configured budget. The same
trained checkpoint is used for both server-side metric evaluation
and on-device measurement.

% =====================================================================
\section{Results}
\label{sec:results}

This section reports the main forecasting accuracy of FEATHer and the
eleven baselines on the eight long-term benchmarks introduced in
Section~\ref{subsec:datasets}, using the protocol specified in
Section~\ref{subsec:protocol} and the evaluation rules in
Section~\ref{subsec:evaluation}.

\subsection{Main Forecasting Accuracy}
\label{subsec:results-main}
Tables~\ref{tab:main-mse} and~\ref{tab:main-mae} report the
five-seed average MSE and MAE for every (model, dataset, horizon)
cell; the best result per row is in \textbf{bold}, the second best
\underline{underlined}. Per-cell standard deviations and the full
$5\!\times\!5$ pairwise Wilcoxon $p$-value matrix are provided in
the appendix. Tables are generated automatically from
\texttt{results/fcst\_results.csv} by
\texttt{tools/paper/main\_table.py}.

\TODO{paste main MSE table from \texttt{python tools/paper/main\_table.py --metric MSE --format latex}}

\TODO{paste main MAE table from \texttt{python tools/paper/main\_table.py --metric MAE --format latex}}

\paragraph{Where FEATHer wins.}
\TODO{summarize per-dataset wins after sweep lands; expected pattern
based on the ultra-light regime: FEATHer competitive on
ETTh\{1,2,m1,m2\}, Weather, Exchange; capacity-rich baselines
(PatchTST, iTransformer, MDMLP-EIA) stronger on
Electricity / Traffic. Report \emph{honestly} per R8 \#1.}

\paragraph{Statistical significance.}
\TODO{paste Wilcoxon table from \texttt{python tools/paper/wilcoxon.py --reference FEATHer}, and call out which pairs are flagged at $p<0.05$ two-sided.}

\subsection{Parameter Efficiency vs Accuracy}
\label{subsec:results-pareto}
\TODO{Pareto plot: x = log10(params) from Table~\ref{tab:params};
y = mean MSE across all 32 cells. FEATHer should sit on the
Pareto frontier alongside SparseTSF / DiPE-Linear at the bottom
of the parameter axis. Generated by
\texttt{tools/paper/main\_table.py --pareto} (to be added).}

% =====================================================================
\section{Ablation Study}
\label{sec:ablation}

Following the absorption of the 30 FEATHer variants into the main
registry, we evaluate every ablation cell with the same five-seed
protocol as Section~\ref{sec:results}; this directly closes
R8 \#6 (``ablations only on $H{=}96$ or $H{=}720$'') by reporting
every variant on \emph{all four} horizons.

\subsection{Multiscale ablation (R5 \#5)}
We disable individual branches and combinations to isolate their
contribution. The 15 multiscale variants
(\texttt{FEATHer\_ms\_P}, \texttt{...\_H}, \texttt{...\_M},
\texttt{...\_L}, \texttt{...\_PH}, \ldots, \texttt{...\_PHML}) cover
every non-empty subset of $\{p, h, m, l\}$.

\TODO{paste table from \texttt{python tools/paper/ablation\_table.py --axis ms --metric MSE --format latex}.}

\subsection{Gating ablation (R2 \#19, \#20)}
Four variants control the gating mechanism:
\texttt{FEATHer\_gate\_\{none, uniform, softmax, fft\}}. The
\texttt{none} variant uses no gating at all (concatenation +
linear), \texttt{uniform} fixes $g_b = 1/B$, \texttt{softmax}
learns logits without spectral input, and \texttt{fft} is the
full FEATHer gate.

\TODO{paste table from \texttt{python tools/paper/ablation\_table.py --axis gate}, then add 2--3 sentences explaining the gate's contribution and addressing R2 \#20 (uniform vs. no gating).}

\subsection{DTK ablation (R5 \#7, R6 \#4b)}
Variants \texttt{FEATHer\_dtk\_\{none, mlp, shallow, full\}}
replace the depthwise temporal kernel with (none) identity,
(mlp) a feed-forward layer, (shallow) a depthwise conv with
$k_{\text{temp}}{=}3$, or (full) the default $k_{\text{temp}}{=}7$.

\TODO{paste table from \texttt{python tools/paper/ablation\_table.py --axis dtk}, then briefly state whether per-branch DTK (R5 \#7) would help or hurt; current code keeps DTK shared (Section~\ref{sec:method}).}

\subsection{Head ablation}
Variants \texttt{FEATHer\_head\_\{linear, mlp, conv, spk\}}
replace the SPK head with simpler alternatives.

\TODO{paste table from \texttt{python tools/paper/ablation\_table.py --axis head}.}

\subsection{Complexity sweep (R8 \#6)}
Variants \texttt{FEATHer\_complexity\_\{half, full, double\}}
scale the latent width $S$ to $S/2$, $S$, $2S$ respectively.

\TODO{paste table from \texttt{python tools/paper/ablation\_table.py --axis complexity}.}

% =====================================================================
\section{Robustness on Sensor Faults}
\label{sec:robustness}

The IoT-J selling point of this work is that FEATHer remains usable
under the four classes of sensor fault that occur in field
deployments (R6 \#7a). We sweep $19\,200$ corruption conditions on
four representative datasets ---
\textbf{ETTh1} (canonical), \textbf{Weather} (real sensors),
\textbf{Electricity} (smart grid), \textbf{SML} (smart home) ---
and reuse the trained Phase 4 checkpoints, so each robustness number
reflects the deployed model exactly.

\subsection{Fault model}
\label{subsec:robust-faults}
We consider four independent fault axes that are documented at
multiple severity levels in \texttt{utils/noise.py}:
\begin{itemize}
\item \textbf{Additive Gaussian noise} with
       $\sigma \in \{0.01, 0.05, 0.10, 0.20, 0.50\}$ on
       the standardized input.
\item \textbf{Missing values} with random masking probability
       $p \in \{0.05, 0.10, 0.20, 0.30, 0.50\}$; missing entries
       are filled with the last-observation-carry-forward rule.
\item \textbf{Impulse outliers} with rate
       $p \in \{0.01, 0.02, 0.05, 0.10\}$ and magnitude $5\sigma$.
\item \textbf{Quantization} to
       $b \in \{8, 6, 5, 4, 3\}$ bits via uniform mid-tread
       quantization of the standardized input.
\end{itemize}
The full grid produces $20$ corruption conditions per dataset
(one clean + $5{+}5{+}4{+}5$ severities); with $12$ models
$\times$ $4$ datasets $\times$ $4$ horizons $\times$ $5$ seeds, the
sweep yields $19\,200$ rows in
\texttt{results/robust\_results.csv}.

\subsection{Aggregate degradation}
\label{subsec:robust-summary}
\TODO{paste table from \texttt{python tools/paper/robust\_summary.py --metric MSE}, which reports mean degradation $\Delta\mathrm{MSE}$ per (model, fault axis) cell averaged across severity levels.}

\subsection{Heatmaps}
\label{subsec:robust-heatmaps}
\TODO{include \texttt{figures/robust\_heatmap\_\{ETTh1,Weather,Electricity,SML\}.pdf} generated by \texttt{python tools/paper/robust\_summary.py --output\_dir manuscript/figures/}.}

\subsection{Discussion}
\TODO{after data lands: identify which fault axis hits which model
family hardest; expected: Transformer family more sensitive to
missing-value bursts, SparseTSF/DiPE-Linear more sensitive to
quantization at $b{\le}4$ bits, FEATHer's spectral gating provides
adaptive resilience to Gaussian noise.}

% =====================================================================
\section{Edge Deployment}
\label{sec:edge}

This section quantifies the on-device cost of every method in
Table~\ref{tab:params}. We target the
\textbf{LM3S6965EVB} (ARM Cortex-M3 at 50\,MHz, $256$\,KB flash,
$64$\,KB on-chip SRAM, $3.3$\,V $V_{\mathrm{DD}}$), and report two
complementary measurements: a fast \emph{simulation-based estimator}
that produces a number for every (model, dataset, horizon, RAM
budget) cell (Section~\ref{subsec:edge-estimator}), and a
\emph{QEMU validation path} that verifies memory fit and functional
correctness of the quantized firmware under emulation
(Section~\ref{subsec:edge-qemu}). All numbers in this section assume
$8$-bit integer weights and activations (post-training quantization
following the CMSIS-NN reference workflow).

\subsection{Simulation-based estimator}
\label{subsec:edge-estimator}
Our estimator inspects each \texttt{nn.Module} via a forward hook,
records the activation byte size and FLOP count of every leaf
operator, and applies the cost model in
Table~\ref{tab:m3-profile}. Peak RAM is computed as a sliding window
over consecutive live activations, mirroring the arena planner used
by TFLite-Micro. Latency is converted via the Cortex-M3
cycles-per-MAC figure published by CMSIS-NN
($3$ cycles/MAC at int8); energy is the product of latency, supply
voltage, and active-mode current.

\begin{table}[t]
\centering
\caption{Cortex-M3 cost profile used by the estimator.}
\label{tab:m3-profile}
\small
\begin{tabular}{lr}
\toprule
Parameter & Value \\
\midrule
Clock frequency        & $50$\,MHz \\
Supply voltage         & $3.3$\,V \\
Active-mode current    & $90$\,mA (CMSIS-NN inference, conservative) \\
Cycles/MAC (int8)      & $3.0$ \\
SRAM capacity          & $64$\,KB \\
Flash capacity         & $256$\,KB \\
Runtime code overhead  & $6$\,KB (CMSIS-NN kernels + main loop) \\
Runtime RAM overhead   & $2$\,KB (stack + bookkeeping) \\
\bottomrule
\end{tabular}
\end{table}

\paragraph{Per-model cost on ETTh1, $H{=}96$.}
Table~\ref{tab:edge-etth1} reports a single representative cell
($D{=}7$, $L{=}96$, $H{=}96$) sorted by parameter count.

\begin{table*}[t]
\centering
\caption{Estimated on-device cost on ETTh1 with $L{=}H{=}96$,
$8$-bit integer inference on Cortex-M3 at $50$\,MHz. The right-most
three columns indicate whether the model fits under
$16$/$32$/$64$\,KB RAM budgets, with the $256$\,KB flash limit
applied jointly.}
\label{tab:edge-etth1}
\small
\begin{tabular}{lrrrrrrccc}
\toprule
Model & Params & Flash & Peak RAM & Arena & Latency & Energy
      & \multicolumn{3}{c}{Deployable at} \\
\cmidrule(lr){8-10}
      &        & (KB)  & (KB)     & (KB)  & (ms)    & (mJ)
      & 16\,KB & 32\,KB & 64\,KB \\
\midrule
SparseTSF      &        41 &    6.0 &    4.6 &  1.3 &    0.31 &   0.09 & \checkmark & \checkmark & \checkmark \\
DiPE-Linear    &       339 &    6.3 &    5.9 &  2.6 &    0.16 &   0.05 & \checkmark & \checkmark & \checkmark \\
\textbf{FEATHer (ours)} &  \textbf{453} &  \textbf{6.4} &  \textbf{6.1} &  \textbf{2.8} &  \textbf{4.03} &  \textbf{1.20} & \checkmark & \checkmark & \checkmark \\
FITS           &     1\,200 &    7.2 &    3.6 &  0.3 &    0.48 &   0.14 & \checkmark & \checkmark & \checkmark \\
DLinear        &    18\,624 &   24.2 &    5.3 &  2.0 &    7.74 &   2.30 & \checkmark & \checkmark & \checkmark \\
TQNet          &    95\,240 &   99.0 &    6.8 &  3.5 &   24.25 &   7.20 & \checkmark & \checkmark & \checkmark \\
TimeMixer      &    75\,497 &   79.7 &   66.3 & 63.0 &  649.4  & 192.9  & --- & --- & --- \\
LMS-AutoTSF    &   181\,318 &  183.1 &   28.6 & 25.3 &  111.0  &  33.0  & --- & \checkmark & \checkmark \\
PatchTST       &   548\,707 &  541.8 &   76.8 & 73.5 & 2083.6  & 618.8  & --- & --- & --- \\
TimesNet       &   605\,479 &  597.3 &   27.3 & 24.0 &   58.3  &  17.3  & --- & --- & --- \\
iTransformer   &   841\,568 &  827.8 &   10.3 &  7.0 &  352.3  & 104.6  & --- & --- & --- \\
MDMLP-EIA      & 1\,671\,446 & 1638.3 &   9.2 &  5.9 &  699.6  & 207.8  & --- & --- & --- \\
\bottomrule
\end{tabular}
\end{table*}

Three points stand out. First, only the six smallest methods
(SparseTSF, DiPE-Linear, FEATHer, FITS, DLinear, TQNet) fit the
strict $16$\,KB RAM budget. Second, the deployability failure mode
varies: TimeMixer and PatchTST are blocked by activation memory
($66.3$ and $76.8$\,KB respectively), while iTransformer and
MDMLP-EIA are blocked by flash ($828$ and $1638$\,KB against the
$256$\,KB chip limit), and TimesNet exceeds the flash budget by
roughly $2.3{\times}$. Third, latency tracks FLOPs as expected, and
FEATHer's $4.03$\,ms response time is comfortably below the
$10$\,ms streaming-inference deadline of typical industrial PLC
loops, while consuming an estimated $1.2$\,mJ per inference.

\paragraph{Deployability across all datasets.}
Aggregating across the eight datasets and four horizons
(32 cells per model), Table~\ref{tab:edge-counts} shows that the
ultra-light methods (FEATHer, SparseTSF, DiPE-Linear, FITS) fit at
least $20$ of $32$ cells even at the $16$\,KB budget, while the
attention-based and MLP-mixing baselines fail in every cell of
every budget. The drop on the $16$\,KB column for DLinear and TQNet
is driven by the high-channel Electricity and Traffic datasets,
where the activation tensors grow with $D$ and exceed the
tightest budget.

\begin{table}[t]
\centering
\caption{Deployability count (out of $32$ cells $=$ $8$ datasets
$\times$ $4$ horizons) by model and RAM budget.}
\label{tab:edge-counts}
\small
\begin{tabular}{lccc}
\toprule
Model               & $16$\,KB & $32$\,KB & $64$\,KB \\
\midrule
FITS                & $23$ & $24$ & $24$ \\
\textbf{FEATHer (ours)} & $\mathbf{22}$ & $\mathbf{23}$ & $\mathbf{24}$ \\
SparseTSF           & $22$ & $23$ & $24$ \\
DiPE-Linear         & $20$ & $23$ & $24$ \\
TQNet               & $19$ & $23$ & $24$ \\
DLinear             & $16$ & $23$ & $24$ \\
LMS-AutoTSF         &  $0$ &  $5$ &  $6$ \\
PatchTST, TimeMixer, TimesNet, &  &  & \\
\quad iTransformer, MDMLP-EIA & $0$ & $0$ & $0$ \\
\bottomrule
\end{tabular}
\end{table}

\subsection{QEMU validation path}
\label{subsec:edge-qemu}
The estimates in Table~\ref{tab:edge-etth1} are first-order: they
assume the CMSIS-NN cycles-per-MAC figure, a fixed runtime overhead,
and an arena planner that approximates the Tensor\-Flow Lite Micro
allocator. To bound the estimator's error, we additionally run the
$8$-bit quantized FEATHer firmware inside QEMU
(\texttt{qemu-system-arm -M lm3s6965evb -kernel} \texttt{a.elf}
\texttt{-nographic -semihosting}). QEMU provides
\emph{instruction-level} (not cycle-accurate) emulation of the
Cortex-M3 core, so we are explicit about what each claim rests on.
The emulation \emph{verifies} three things directly:
(i)~the quantized firmware executes the full inference correctly,
producing outputs that match the int8 PyTorch reference;
(ii)~the static memory fit is real --- the linker script enforces
the $64$\,KB SRAM / $256$\,KB flash limits of the target, so a
model that exceeds either budget fails to build at all; and
(iii)~the deterministic instruction count of one inference, read
back via QEMU's \texttt{icount} mechanism over the semi-hosting
channel. Latency is then \emph{derived} from the verified
instruction count through the cycles-per-instruction model of
Table~\ref{tab:m3-profile}, anchored to the CMSIS-NN
cycles-per-MAC figure. The firmware is cross-compiled with
\texttt{arm-none-eabi-gcc}~12.2.1
(\texttt{-Os -mthumb -mcpu=cortex-m3}) and links against CMSIS-NN.
The full validation pipeline --- model quantization, C-array
codegen, firmware build, QEMU run, and result aggregation --- is
released alongside the code under \texttt{deployment/cortex\_m3/}.

\emph{Reporting policy.} We separate \emph{verified} from
\emph{modeled} quantities. Memory fit (the deploy columns) and
functional correctness are verified: the linker and the emulated
run make them hard pass/fail outcomes. Latency is modeled from the
verified instruction count via the cost profile of
Table~\ref{tab:m3-profile}, and energy is modeled as latency
$\times$ supply power; both are labeled estimates wherever they
appear. Future work includes a measurement on a physical Cortex-M3
board (e.g., STM32F103) with energy instrumentation; we expect the
on-board figures to track the estimator within the envelope implied
by the stated cost model.

% =====================================================================
\section{Conclusion}
\label{sec:conclusion}

We presented FEATHer, a deployable forecasting system designed for
long-horizon prediction on the kilobyte-scale weight budgets of
industrial edge controllers. FEATHer combines four lightweight
components --- multiscale temporal decomposition, a shared Dense
Temporal Kernel, frequency-aware gating, and a parameter-minimal
Sparse-Period Kernel --- inside a compact architecture that holds a
sub-1K parameter budget across edge-typical channel counts.

\TODO{one-paragraph quantitative summary once Phase 4 lands.
Expected shape: ``FEATHer reaches an average MSE of X across
eight datasets and four horizons with $\le 1\,000$ parameters,
matching or surpassing baselines with $10^{3}$--$10^{5}$
parameters in the ultra-light regime, while remaining
$\ge K$\,KB below the $16$\,KB Cortex-M3 RAM budget on every
configuration that satisfies the sub-1K target.''}

\paragraph{Limitations.}
Three are inherent to the design choices made here and we record
them honestly:
(i)~On high-channel datasets (Traffic with $D{=}862$, Electricity
with $D{=}321$), capacity-rich Transformer baselines retain a
measurable accuracy advantage that the ultra-light regime cannot
close; this is the deliberate trade-off behind the sub-1K
parameter budget.
(ii)~The ``sub-1K parameter'' framing is
\emph{configuration-dependent} (Section~\ref{subsec:edge-estimator})
--- the count holds for small-$D$ datasets but grows on Weather and
Solar; we report this explicitly rather than averaging over
configurations.
(iii)~The Cortex-M3 deployability numbers in
Section~\ref{sec:edge} are simulation-based estimates. The QEMU
validation path described in Section~\ref{subsec:edge-qemu} narrows
the gap to a cycle-accurate run, but a physical-board measurement
campaign with energy instrumentation remains future work.

\paragraph{Future directions.}
We see three productive extensions.
First, replacing the validation-set search over $P$ and the gate
spectrum bin count with a mechanism that adapts them
\emph{automatically} to a dataset's dominant period --- without
sacrificing the deployability guarantee. Second, multimodal sensor fusion: combining the
DTK with a lightweight cross-sensor mixer would let FEATHer
ingest heterogeneous industrial signals (vibration, temperature,
voltage) without abandoning the sub-1K target. Third, hardware/
loss co-optimization: the spectral-separation auxiliary loss
$\mathcal{L}_{\mathrm{spec}}$ leaves a hook for explicit
energy-aware training, which is well aligned with the IoT-J
mission of low-power on-device inference.

The full code (including the Phase 4 sweep orchestrator,
the Cortex-M3 estimator, and the manuscript build) is released
to make every number in this paper reproducible.

% =====================================================================
\bibliographystyle{IEEEtran}
\bibliography{feather}

\end{document}
